\subsection*{一个字符串在什么范围内唯一}

客户端 A 与客户端 B 都可能发来 K1。假如 A 属于账户 1，B 属于账户 2，服务只用键字符串查询，账户 2 就可能命中账户 1 的购买。这既会误合并购买，也可能泄露订单与账户信息。这种确定一个标识在哪个范围内解释的规则，称为命名空间；空间名称与键共同决定查找对象。服务器需要先完成认证，解析调用者被允许操作的账户，再构造查找身份。

此时先使用组合身份
\[
 S=(\mathrm{account\_id},\,\mathrm{action},\,\mathrm{idempotency\_key}),
\]
其中 \code{action} 固定为 \code{purchase.create.v1}，表示第一版创建购买动作。账户 1 的 K1 对应 \code{(1, purchase.create.v1, K1)}，账户 2 的 K1 对应 \code{(2, purchase.create.v1, K1)}。同一账户的退款动作进入 \code{refund.create.v1} 空间。字符串相同而账户或动作不同，组合身份仍然不同。后面为身份增加有效期时，还会加入登记会话这一维。

\code{account\_id} 来自已验证的授权上下文。客户端把请求体里的 \code{account\_id} 改成 2，并不能令服务进入账户 2 的空间。租户是使用同一套服务但数据与权限各自隔离的一组用户或组织。面向企业租户的系统可以先用 \code{tenant\_id}（租户编号）划分，再在租户内加入账户或操作主体；具体选择取决于一个业务操作允许影响的账户范围。若一笔操作可以同时影响多个账户，应由已认证租户和明确业务主体界定它的所有权，不能随意选一条资金记录的账户代替操作空间。

\begin{figure}[htbp]
\centering
\begin{tikzpicture}[x=1cm,y=1cm]
\node[booknode,text width=4.7cm] at (2.6,0) {账户 1，创建购买\\(1, purchase.create.v1, K1)\\对应 P1};
\node[booknode,text width=4.7cm] at (9,0) {账户 2，创建购买\\(2, purchase.create.v1, K1)\\对应 P8};
\node[booknode,text width=4.7cm] at (2.6,-2) {账户 1，申请退款\\(1, refund.create.v1, K1)\\对应退款 F1};
\node[bookresult,text width=4.7cm] at (9,-2) {只查字符串 K1\\三个不同对象被混在一起\\可能误合并或泄露数据};
\node[booknote,text width=12cm] at (5.8,-3.4)
 {先认证并确定操作范围，再构造完整身份；相同字符串可以存在于不同命名空间};
\end{tikzpicture}
\caption{账户和动作共同确定键的解释范围。加入登记会话后，查找还须携带会话维度。}
\label{fig:ch03-namespaces}
\end{figure}


命名空间也确定了唯一性要求的作用范围。在一个空间内，不同操作应当取得不同键；空间之间的键可以相同。如果账户 1 与账户 2 的 K1 都能独立受理，系统仍符合这项规则。若数据库只在键字符串上加全局唯一约束，账户 2 的正常请求就会遇到账户 1 的记录。反过来，若查询按账户和动作定位，写入时却只检查局部进程中的键，不同服务进程又可能各自受理同一完整身份。查找、唯一检查和恢复都要使用同一个完整对象。

命名空间还要贯穿查询入口。查询 K1 时，服务同样先认证账户，并以完整组合身份读取记录。知道 K1 或 P1 只提供查找线索，访问权限仍由认证与授权决定。随机键难以猜测可以降低偶然枚举的机会，却不等于持有人获得订单访问权。键也应限制长度并按受控方式记录日志，减少任意输入占用存储或污染诊断信息。

API 的路径与动作不总是一一对应。移动端的 \code{/purchases} 与网页端的 \code{/checkout} 如果都是同一购买动作，应在服务内部映射到同一 \code{action}，才能让跨入口重试恢复原购买。若两个入口生成不同空间，相同 K1 会被分别受理。相反，把购买与退款都映射为同一 \code{action}，会让退款命中原购买。动作名表达业务效果与输入解释规则。

接口升级也需要明确安排。若 v2 只增加可选展示字段，可以保留相同动作与指纹规则，旧操作继续按其原版本解释。指纹是由规定输入计算出的摘要，用于快速核对同身份请求是否仍携带原输入，下一节会具体构造。若 v2 引入数量或新的计价规则，使效果发生变化，可以使用 \code{purchase.create.v2}，并规定旧操作仍由 v1 恢复入口处理。客户端在一次待确认购买中途切到新动作空间，相当于提出新操作，升级代码必须保存旧记录的动作版本。

\subsection*{键沿着哪条业务链路传播}

在本书四表服务里，购买、订单和扣款都由一个事务完成。入口可以把 K1 交给购买处理层，处理层建立 K1 到 P1 的关联，订单和扣款继续引用 P1。每一层应显式传递账户与动作，防止底层工具只收到一个脱离范围的字符串。入口日志还应保留本次 R1，方便把一次失败尝试与业务身份并列查询。

若业务拆成购买服务和资金服务，调用关系会多一层。购买服务针对同一 P1 请求一次扣款，扣款重试应继续使用一个稳定的派生身份。例如，可以组合以下四个字段：
\begin{BookCode}
(purchase.create.v1, P1, debit, 1)
\end{BookCode}
其中 \code{debit} 表示子动作，1 表示本次购买中这项扣款的序号。资金服务再加上自己的已认证调用方和扣款动作空间。它收到的重试身份需要由购买服务保存，不能每次网络调用临时生成。

\begin{figure}[htbp]
\centering
\begin{tikzpicture}[x=1cm,y=1cm]
\node[booknode,text width=5.4cm] (p) at (6,0) {父操作：账户、购买动作、K1\\服务保存稳定购买编号 P1};
\node[booknode,text width=3.5cm] (d) at (1.5,-2) {扣款子动作\\(P1, debit, 1)\\扣款重试继续此身份};
\node[booknode,text width=3.5cm] (n) at (6,-2) {通知子动作\\(P1, notice, 1)\\通知重试继续此身份};
\node[booknode,text width=3.5cm] (i) at (10.5,-2) {发票子动作\\(P1, invoice, 1)\\发票重试继续此身份};
\draw[bookarrow] (p)--(d);
\draw[bookarrow] (p)--(n);
\draw[bookarrow] (p)--(i);
\node[booknote,text width=12cm] at (6,-3.6)
 {接收服务再加自己的已认证调用方与动作范围；箭头说明来源和调用关系，跨服务共同提交另需协议};
\end{tikzpicture}
\caption{父操作相同，子动作仍需分别识别；同一子动作重试继续沿用自己的身份。}
\label{fig:ch03-derived-identity}
\end{figure}


同一购买还可能产生通知、发票或积分动作。它们引用 P1 说明业务来源，各自的效果和恢复规则则不同。若直接把 K1 作为资金服务中所有动作的唯一键，第二项合法子动作可能命中第一项；若每次调用都新建子键，超时又会重复受理。派生规则要包含接收方所需的稳定父身份、动作和子项序号，并在重试间保持一致。

可以把派生身份写为 \(S_{\mathrm{child}}=F(S_{\mathrm{parent}},a,i)\)。\(S_{\mathrm{parent}}\) 指父操作的完整身份，\(a\) 是子动作，\(i\) 是同类子动作的稳定序号，\(F\) 是双方约定的组合编码。相同三个输入需要得到相同子身份，已区分的子项需要得到不同子身份。编码应保留字段边界，例如用结构化字段或长度前缀，避免把两组不同字段直接拼成同一个字符串。

这个公式还要求 \(i\) 来自已经确定的业务子项。进程重启后若按另一顺序重新遍历子项，第一项、第二项交换位置，即使 \(F\) 完全确定，身份仍可能指错对象。可以保存子项及其序号，也可以直接使用稳定的子项记录编号。输入来源的稳定性，与编码计算的确定性共同构成派生规则。

退款应创建新的退款操作，并引用原扣款 D1。它改变另一笔资金事实，用户可以对同一原扣款分次申请退款；原购买 K1 无法同时表示这些退款。退款接口要核对累计返还额度及各次退款身份，那是退款自己的业务规则。原购买与后续动作通过引用相连，查询时仍能分别得到购买原结果和每项后续处理结果。

派生身份定义请求对应关系，跨服务更新的共同提交还需要额外机制。购买服务已提交、资金服务结果未知，或扣款已完成、购买记录尚未更新，都需要消息、事务或恢复设计。图中的箭头表示调用与引用；资金和订单的一致性仍要由另外的执行协议维护。
\FloatBarrier
